\section{Trygonometria hiperboliczna} % lang-pl
\label{sekcja_trygonometria}

Lambert wprowadzi do matematyki funkcje hiperboliczne i obliczy pole trójkąta hiperbolicznego (sekcja \ref{sekcja_lambert}), a Taurinus opublikuje w latach 1825--1826 wyniki trygonometrii hiperbolicznej (sekcja \ref{sekcja_schweikart_taurinus}). % lang-pl
W tej sekcji zobaczymy, jak Hartshorne wyprowadza ją bez liczb rzeczywistych i bez funkcji sinh, cosh, tanh, jak wyglądają jej wzory i dwa zastosowania: pole koła oraz trójkąty stowarzyszone Engela. % lang-pl

Hartshorne \cite[s. 403--404]{hartshorne_2000} tłumaczy, że zwykła trygonometria to związki między bokami i kątami trójkąta prostokątnego wraz z regułami rachunku, które pozwalają wyliczyć pozostałe elementy z dowolnych dwóch. % lang-pl
W geometrii hiperbolicznej jest to nawet trochę mocniejsze, bo nie ma trójkątów podobnych, więc dwa kąty $\alpha, \beta$ wyznaczają trójkąt jednoznacznie; komplikuje sprawę tylko to, że $\alpha + \beta < \mathrm{RA}$, więc $\beta$ nie jest wyznaczone przez $\alpha$. % lang-pl
Zamiast logarytmów używa się multiplikatywnej funkcji odległości $\mu$ z sekcji \ref{sekcja_ciala_koncow}, a ważną stałą $k$, która w niektórych książkach zależy od wyboru podstawy logarytmu, w ogóle się nie pojawia: dwie płaszczyzny hiperboliczne nad tym samym ciałem są izomorficzne \cite[s. 296, 404]{hartshorne_2000}. % lang-pl

\begin{definition}
[sinus i cosinus kąta] % lang-pl
	Dla kąta $\alpha$ i $t = \tan\tfrac{\alpha}{2}$ (w sensie sekcji \ref{sekcja_ciala_koncow}) kładziemy $\sin\alpha = \dfrac{2t}{1+t^2}$ i $\cos\alpha = \dfrac{1-t^2}{1+t^2}$. % lang-pl
\end{definition}

Funkcje te spełniają zwykłe tożsamości: $\tan\alpha = \sin\alpha/\cos\alpha$ i $\sin^2\alpha + \cos^2\alpha = 1$ \cite[s. 405--406]{hartshorne_2000} (prop. 42.1). % lang-pl

\begin{theorem}
	W trójkącie prostokątnym o kątach $\alpha, \beta$ przy wierzchołkach $A, B$, i o kątach równoległości $\bar a, \bar b, \bar c$ boków naprzeciw $A, B, C$ zachodzi: % lang-pl
	$\tan\alpha = \cos\bar a\,\tan\bar b$, $\quad \cos\bar b = \cos\alpha\,\cos\bar c$, $\quad \sin\alpha = \cos\beta\,\sin\bar b$; % lang-pl
	$\tan\bar c = \sin\alpha\,\tan\bar a$, $\quad \sin\bar c = \tan\alpha\,\tan\beta$, $\quad \sin\bar c = \sin\bar a\,\sin\bar b$. % lang-pl
\end{theorem}

Dowody: Hartshorne \cite[s. 406--407]{hartshorne_2000} (prop. 42.2, 42.3); wyprowadzenie zaczyna od trójkąta w szczególnym położeniu w układzie końców i wzoru Bolyaia $\tan\tfrac{\bar a}{2} = a^{-1}$ \cite[s. 404--405]{hartshorne_2000}. % lang-pl
Jak zauważa Hartshorne, te wzory i cztery dalsze uzyskane przez zamianę ról wierzchołków dają po jednym równaniu dla każdego podzbioru trzech spośród pięciu wielkości $a, b, c, \alpha, \beta$, więc znając dowolne dwie, można obliczyć pozostałe \cite[s. 407]{hartshorne_2000}. % lang-pl

W pierwszym z zastosowań Hartshorne oblicza zależność między bokiem a kątem trójkąta równobocznego (przykład 42.3.2) i pole koła. % lang-pl
Jeśli ciało końców $F$ zawiera się w ciele liczb rzeczywistych, to pole koła o promieniu, którego multiplikatywna długość wynosi $r$, jest granicą pól wielokątów foremnych wpisanych i równa się $A = \pi(r-1)^2/r$ \cite[s. 407--409]{hartshorne_2000} (prop. 42.4); dla zwykłej długości promienia $x = \ln r$ wychodzi $A = 4\pi\sinh^2(x/2)$, więc pole rośnie szybciej niż $\pi x^2$ (zad. 42.16) \cite[s. 413]{hartshorne_2000}. % lang-pl

Hartshorne \cite[s. 409]{hartshorne_2000} używa tego wzoru do pytania o „kwadraturę koła'' w geometrii nieeuklidesowej: w geometrii euklidesowej nad ciałem liczb konstruowalnych odpowiedź brzmi zawsze „nie'', bo $\pi$ jest przestępna, ale w płaszczyźnie hiperbolicznej odpowiedź bywa raz „tak'', raz „nie'' -- zależnie od promienia. % lang-pl
Jeśli $x = (r-1)^2/r$ jest dyadyczną liczbą wymierną, to $\pi x$ jest sumą miar kątów, którą da się dowolnie połowić, i koło jest równoważne figurze prostoliniowej; dla promienia $r = \tfrac{1}{14}(15+\sqrt{29})$ nad ciałem liczb konstruowalnych już nie, bo kąt $\pi/7$ nie istnieje w takiej płaszczyźnie \cite[s. 409]{hartshorne_2000}. % lang-pl

Drugie zastosowanie to \emph{trójkąty stowarzyszone Engela}: każdemu trójkątowi prostokątnemu o elementach $(a, b, c, \alpha, \beta)$ odpowiada trójkąt prostokątny o elementach $(b, \beta', \bar\alpha', \bar a', \bar c)$ (gdzie prim oznacza kąt dopełniający do $\mathrm{RA}$) \cite[s. 409--410]{hartshorne_2000} (prop. 42.5). % lang-pl
Powtórzenie tej operacji pięć razy daje z powrotem trójkąt wyjściowy, więc jeden trójkąt rodzi cykl pięciu trójkątów stowarzyszonych \cite[s. 411]{hartshorne_2000}; konstrukcję geometryczną daje konstrukcja Bolyaia z sekcji \ref{sekcja_ciala_koncow} (zad. 42.23) \cite[s. 411, 414]{hartshorne_2000}. % lang-pl
\index[persons]{Engel, Friedrich}%

Zadania: prawo sinusów (zad. 42.6) i wzór Łobaczewskiego, analogiczny do prawa cosinusów (zad. 42.7c) \cite[s. 411--412]{hartshorne_2000}; tłumaczenie na funkcje sinh, cosh, tanh (zad. 42.15) \cite[s. 413]{hartshorne_2000}; „trójki pitagorejskie'' w geometrii nieeuklidesowej (zad. 42.25) \cite[s. 415]{hartshorne_2000}. % lang-pl
% Źródła: Hartshorne s. 415--434 (§43); przeczytano dokładnie s. 415--422 i fragmenty s. 424--428.

